\section{Introduction} \label{intro}

Adaptive evolution depends on both the generation of variation and the action of natural selection. These processes are commonly described using two conceptual tools, the first of which is the genotype–phenotype map, characterizing how genetic variation is translated into phenotypic variation \parencite{houle2010phenomics}. 
The second is the adaptive landscape, which determines the relationship between phenotype and fitness \parencite{arnold2001adaptive, steppan2002comparative, Schluter2024-wd}. 
Together, they define how populations can move through spaces of genotypes and phenotypes by determining which phenotypic variants are produced by mutation and which are favored by selection.

Both the adaptive landscape and the genotype-phenotype map are globally complex structures, defined over the space of possible genotypes and phenotypes. Adaptive landscapes can be rugged and multi-peaked \parencite{Martin2013-ez, Martin2020-kn} and may shift seasonally in unpredictable ways \parencite{Stroud2023-nn}. The diversity of forms observed in adaptive radiations illustrates how multiple phenotypic optima can be reached from a single ancestor. Although precise empirical estimates of the ruggedness of adaptive landscapes in nature are scarce, simple single-peaked, static adaptive landscapes may be the exception rather than the rule \parencite{Uyeda2024-ds}. The genotype-phenotype map is no less complex. Phenotypic traits are the outcome of development, a dynamic process through which interactions at the genetic, cellular and tissue levels produce phenotypes \parencite{alberch1991genes, houle2010phenomics, milocco2022evolution}. The genotype–phenotype map therefore varies across sequence space, such that the size and direction of a mutation's effect depend on the genetic background in which it occurs.

Despite this global complexity, populations experience both the adaptive landscape and the genotype-phenotype map only locally at any given point in time. Quantitative genetics formalizes this local perspective, providing tractable representations of both \parencite{lynch1998genetics, melo2016modularity}. The adaptive landscape is locally summarized by the selection gradient, the partial regression of trait values on relative fitness, capturing the direction and strength of directional selection on phenotypic traits \parencite{lande1979quantitative, Lande1983-lf}. The G-matrix encodes the additive genetic variances and covariances among traits present in a population at a given point in time and describes the directions in phenotype space along which heritable variation is available \parencite{Schluter1996-uo, hansen2008measuring}. Together, these quantities predict short-term evolutionary change (but see  \cite{milocco2020evolution, rice2002general}). 
As a statistical measure of standing genetic variation, the G-matrix is shaped both by processes that introduce and modify variation, such as mutation and recombination, and by processes that remove it, including selection and drift \parencite{jones2004evolution, melo2015directional}. The input from mutation is captured by the M-matrix, the variance-covariance matrix of phenotypic variation introduced by mutation each generation \parencite{jones2007mutation, Jones2014-wj}. Because the M-matrix describes how mutations arising at the level of the genotype are expressed at the level of the phenotype, it can be interpreted as a local summary of the genotype-phenotype map and, given unbiased mutation at the genotypic level, as a local measure of developmental bias.

The strength of this local framework is also its limitation. When the adaptive landscape or the genotype-phenotype map are globally complex, local representations fail to capture properties that matter over longer timescales. A multi-peaked adaptive landscape means that the selection gradient, defined locally around the current population mean, cannot determine which peak a population will ultimately reach \parencite{melo2016modularity}. Similarly, the G-matrix is not fixed and can evolve as selection gradients change in different regions of the adaptive landscape and as the M-matrix changes over different parts of the genotype-phenotype map \parencite{ cai2025disentangling, milocco2022evolution, rice2002general}. When complex genotype-phenotype maps and adaptive landscapes interact, the resulting long-term evolutionary dynamics may be fundamentally difficult to anticipate from local quantities alone. 

We here address three key questions regarding the relationship of genotype-phenotype map and adaptive landscape. First, how do different genotype-phenotype maps affect evolutionary trajectories and outcomes? Second, how does the interplay between adaptive landscape complexity and genotype-phenotype map affect our ability to predict evolution? And third, what are the consequences for long-term evolution if developmental bias can evolve?

We tackle these questions by independently varying the complexity of both the genotype-phenotype map and the adaptive landscape. We proceed in three stages of increasing complexity. First, we consider the scenario where the genotype-phenotype map is linear over phenotype space, and the adaptive landscape is characterized by a single peak. Second, we increase the complexity of the adaptive landscape by introducing multiple fitness optima. Even with a simple genotype-phenotype map, this creates qualitatively new dynamics: populations with identical mean phenotypes but different M-matrices can evolve to different optima, demonstrating that genetic variation can determine evolutionary outcomes on a complex landscape. Third, we allow the genotype-phenotype map itself to become nonlinear, enabling the M-matrix to evolve. This introduces a further layer of complexity: populations that begin with very similar patterns of genetic covariation can diverge dramatically over time as their variance structures evolve differently, ultimately reaching different optima even when starting from very similar conditions.

